Write $\overline X$ for the projective cone, with vertex $P=(0:\cdots:0:1)$.

(a) On $U_i=V\cap D_+(x_i)$, the coordinate $x_{n+1}/x_i$ identifies $\pi^{-1}(U_i)$ with $U_i\times\mathbb A^1$. Hence $\overline X-P$ satisfies $(*)$. Its prime divisors are either inverse images of prime divisors of $V$, or dominate $V$. The latter restrict to prime divisors of $\mathbb A^1_{K(V)}$, so subtracting divisors of irreducible polynomials in its coordinate leaves only inverse images of divisors of $V$. Thus $\pi^*$ is surjective. If $\pi^*D=(f)$, then $f$ has neither zeros nor poles on the generic affine line, so $f\in K(V)^*$ and $D=(f)$, by the valuations along inverse images. Thus $\pi^*$ is injective. Since $\operatorname{codim}_{\overline X}P\geq2$, (6.5) gives $\operatorname{Cl}\overline X\cong\operatorname{Cl}(\overline X-P)\cong\operatorname{Cl}V$.

(b) If $H=V_+(h)\subseteq\mathbb P^n$, then on $\overline X-P$,
\[(x_{n+1}/h)=V-\pi^*(V.H).\]
Consequently $[V]=\pi^*[V.H]$. Applying (6.5) to $\overline X-X=V$, and then removing the vertex, gives
\[0\to\mathbb Z\xrightarrow{1\mapsto[V.H]}\operatorname{Cl}V\to\operatorname{Cl}X\to0.\]
The first map is injective because $\deg(V.H)=\deg V>0$ by \exref{II.6.2}.

(c) The ring $S(V)$ is normal exactly when $V$ is projectively normal. By (6.2), it is a UFD exactly when it is normal and $\operatorname{Cl}X=0$. By (b), the latter condition means that $\operatorname{Cl}V$ is generated by $[V.H]$; this generator has infinite order.

(d) Restriction to $\operatorname{Spec}\mathcal O_P$ is surjective, with kernel generated by prime divisors $Y$ avoiding $P$. Let $\mathfrak p$ define such a divisor, and choose $h\in\mathfrak p$ with constant term $1$. On the generic punctured line over $V$, $\mathfrak p$ is generated by an irreducible polynomial $f\in K(V)[t]$, normalized by $f(0)=1$. Write $h=fg$ with $g\in K(V)[t]$ and $g(0)=1$. The horizontal part of $(f)$ is precisely $Y$.

For a vertical prime divisor $Z$, choose a cone coordinate trivializing the punctured line near its generic point. Its valuation on a polynomial is the minimum of the valuations of the coefficients. Since the constant terms of $f,g$ are $1$, both valuations are nonpositive. Since $h$ has constant term $1$, it belongs to no homogeneous prime, so $v_Z(h)=0$. Hence $v_Z(f)=v_Z(g)=0$. This applies in every trivialization, whose coordinate changes only rescale $t$. Thus $(f)=Y$, and the kernel is zero.
